\problemname{Pyramid Dice}

In this problem, we use three tetrahedra, four-sided solids, as dice. A thrown tetrahedral die is read by lifting it up and checking the number on the bottom. The twelve sides in total each have a unique number, $1...12$. We are given a series of 20 throws with the three dice and must use these to determine which numbers are on the same die. The dice are read in arbitrary order after each throw.

\section*{Input}
The input consists of 20 dice throws. Each throw is on a separate line and consists of three numbers between $1$ and $12$. The input is designed so that there is exactly one solution.

\section*{Output}
The output should consist of one line for each die. Each line should contain four integers --- the numbers on the die. You must output the numbers in sorted order. You should first output the die that has the number 1. Among the two remaining dice, you should first output the one with the smallest minimum number.
